
\documentclass[../../../physics-standard_questions_all.tex]{subfiles}


\begin{document}

\setlength{\abovedisplayskip}{2mm}
\setlength{\belowdisplayskip}{1mm}

一様な金属導線内での自由電子（質量$m$，電荷$-e$）の運動について考える．
導線内では，自由電子は，電場による加速と金属イオンとの衝突による減速をくり返しながら移動する．
導線に沿って一様な電場（強さ$E$）を加えたとき，
加速された自由電子は一定時間$r$ごとに金属イオンと衝突をくり返し，
衝突のたびに速さがゼロになると考えることにする．

\begin{enumerate}
    \item 自由電子の平均の速さ$\overline{v}$を求めよ．ただし，$\overline{v}$は，
        金属イオンとの衝突直前の速さの$\displaystyle \frac{1}{2}$倍とする．

    \item 導線を流れる電流$I$を求めよ．
        ただし，自由電子の数密度（単位体積あたりの個数）を$n$とし，導線の断面積を$S$とする．

    \item 導線の長さ$L$の部分に注目し，その両端の電位差（電圧）の大きさを$V$とする．
        $I$は$V$に比例することを示し，この導線の抵抗率$\rho$を求めよ．

    \item 注目部分に含まれる自由電子が単位時間あたりに衝突によって失う
        運動エネルギーの総量$P$を求めよ．ただし，$I$，$V$のみを用いて表せ．
    
    
\end{enumerate}

\end{document}



